\documentclass[10pt,a4paper]{amsart}
\usepackage[margin=19mm]{geometry}
\usepackage{amsmath,amssymb,mathtools}
\usepackage[scr=rsfs,cal=euler]{mathalfa}
\usepackage{xfrac}
\usepackage[unicode,hidelinks,bookmarks=false]{hyperref}
\newcommand{\ie}{i.e.\ }
\newcommand{\cf}{cf.\ }
\newcommand{\resp}{respectively\ }
\newcommand{\Subsection}{\subsection}
\providecommand{\nobreakdash}{\nobreak\hbox{-}}
\setlength{\emergencystretch}{2em}
\multlinegap0pt
\usepackage{bm}
\usepackage{bbm}
\usepackage{dsfont}
\usepackage{xspace}
\usepackage{cancel}
\usepackage{mathtools}
\usepackage{amsthm}
\makeatletter
\@ifundefined{theorem}{\newtheorem{theorem}{Theorem}}{}
\@ifundefined{lemma}{\newtheorem{lemma}[theorem]{Lemma}}{}
\makeatother
\makeatletter
\DeclareMathOperator{\SPD}{\sf{SPD}}
\newcommand{\bbZ}{\mathbb{Z}}
\newcommand{\bfP}{\mathbf{P}}
\expandafter\ifx\csname shift\endcsname\relax
\newenvironment{shift}[1][(0,0)]{
    \begin{scope}[shift={#1}]
}{
    \end{scope}
}
\fi
\makeatother
\title[Cauchy identities for Grothendieck polynomials and a dual RS]{Cauchy identities for Grothendieck polynomials and a dual RSK correspondence through pipe dreams}
\author{Hugh Dennin}
\date{Source: arXiv:2506.21052v1 (CC BY 4.0)}
\begin{document}
\maketitle

\noindent\emph{Source and licence.} Cauchy identities for Grothendieck polynomials and a dual RSK correspondence through pipe dreams. Version arXiv:2506.21052v1 (CC BY 4.0), distributed under the Creative Commons Attribution 4.0 International licence, as recorded in the arXiv licence metadata for this exact version and shown on the abstract page https://arxiv.org/abs/2506.21052v1. Full rights notice: licenses/arxiv-2506.21052-CC-BY-4.0.txt.\par\smallskip

\noindent\textit{Editorial scope.} Counted unit: the Lemma whose source label is \texttt{lem:SW corner of big ladder} in the article (source0.tex lines 1621--1637, source label \texttt{lem:SW corner of big ladder}), delivered together with the article's own complete original proof of it (source0.tex lines 1639--1674, one \texttt{proof} environment of 36 lines). No mathematics is written, repaired or re-derived: statement and proof are the article's own text line for line. Required context retained uncounted: none. Every definition, display and label of those lines is retained unchanged. Omitted from this package and not redistributed: 1--1620, 1638--1638, 1675--3527. Nothing inside a retained excerpt is omitted or altered: every retained body is delivered exactly as the article's lines are written, so no omission receipt of any kind accompanies this package. Citations: the retained excerpts carry no citation command at all, so no bibliography entry of the article is transcribed and no reference is redistributed. Reference adaptation: none was needed and none was made; every reference inside the delivered unit resolves inside it, so no unresolved reference or citation is produced. Printed slips kept literal: no printed word, symbol, hypothesis or number of the counted statement or of its proof was altered. Layout and class: the article's own class is \texttt{amsart} with packages including amsfonts, amsmath, amssymb, amsthm, commath, eucal, mathrsfs, mathtools, nccmath, tikz, tikz-cd; the wrapper is the lane's pinned \texttt{amsart} editorial wrapper at 10pt on A4, which restates the theorem-like environments the retained text needs and adds the article's own macro definitions that the retained text needs (\texttt{\textbackslash SPD}, \texttt{\textbackslash bbZ}, \texttt{\textbackslash bfP}), with extra wrapper packages (bm, bbm, dsfont, xspace, cancel, mathtools). The delivered numbering is the unit's own. Assets: no retained span contains a figure environment, a graphics-inclusion command, a PDF-page inclusion, a table or a TikZ picture; no image file is copied into this package, the package declares no asset dependency and no image file is redistributed with it. Licence: version arXiv:2506.21052v1 of this article is distributed under the Creative Commons Attribution 4.0 International licence, as recorded in the arXiv licence metadata for this exact version and shown on the abstract page https://arxiv.org/abs/2506.21052v1. A full rights notice is delivered at licenses/arxiv-2506.21052-CC-BY-4.0.txt. Characters: every retained excerpt is pure ASCII, so the delivered engine, XeLaTeX with its default Latin Modern text font, reports no missing character.
\begin{lemma}
\label{lem:SW corner of big ladder}
    Let $\bfP\in \SPD(w)$.
    Suppose $(i,j)$ is the SW corner of a big ladder $L$ during step $j$ of $Y^+\bfP$.
    Let $(i_0,j)$ (resp. $(i,j_0)$) denote the first tile north (resp. east) of $(i,j)$ with no checkers in $\bfP$ (so $(i_0,j+1)$ is the NE corner of $L$).
    
    \begin{itemize}
    \item[(i)]
    For each $j<j'< j_0$, $(i,j')$ has no checkers in $Y^+_{\geq j'} \bfP$ and, hence, appears in a big ladder $L'$ during step $j'$ of $Y^+\bfP$.
    
    \item[(ii)]
    If $(i,j)$ appears in a big chute $C$ during step $i$ of $X^+\bfP$, then $(i,j)$ is the NE corner of $C$.
    
    \item[(iii)]
    For each $i_0 < i'\leq i$, $(i',j)$ has a red checker in $X^+_{\geq i'} \bfP$.
    \end{itemize}
\end{lemma}

\begin{proof}
We first show (i) by induction on $j< j'< j_0$.
If $j' = j+1$, then $(i,j+1)$ has no checkers in $Y^+_{\geq j+1} \bfP$ since it is the SE corner of the big ladder $L$ during step $j$ of $Y^+\bfP$.
If $j' > j+1$, then by induction $(i,j'-1)$ has no checkers in $Y^+_{\geq j'-1}$ and appears in a big ladder $L'$ during this step; these conditions imply that $(i,j')$ must not have any checkers in $Y^+_{\geq j'}\bfP$.
In any case, $(i,j')$ has checkers in $\bfP$ but not in $Y^+\bfP$, so it must appear in big ladder during step $j+1$ of $Y^+\bfP$.
This completes (i).

We will next show that (iii) is a consequence of (ii), so assume that (ii) holds.
We proceed by downwards induction on $i\geq i'> i_0$.
If $i' = i$, notice that $(i,j)$ being the SW corner of a big ladder $L$ during step $j$ of $Y^+\bfP$ means that $(i,j)$ has a red checker in $\bfP$.
Hence the only way for $(i,j)$ to not contain a red checker in $X^+_{\geq i}$ is if it is contained in a middle column of a big chute during step $i$ of $X^+\bfP$, but this cannot happen by (ii).

Now, let $i_0<i'<i$ and assume there is a red checker at $(i'+1,j)$ in $X^+_{\geq i'+1}\bfP$.
There is also is a checker at $(i',j)$ in $\bfP$ and, hence, in $X^+_{\geq i'+1}\bfP$ since $L$ is a big ladder.
If $(i',j)$ is contained in a chute during step $i'$ of $X^+\bfP$, then the checkers at $(i',j)$ and $(i'+1,j)$ in $X^+_{\geq i'+1}\bfP$ get swapped in $X^+_{\geq i'}\bfP$, so $(i',j)$ has a red checker in $X^+_{\geq i'}\bfP$.
Otherwise $(i',j)$ is not contained in a chute during step $i'$ of $X^+\bfP$, in which case $(i',j)$ must already have a red checker in $\bfP$ and will continue to have this red checker in $X^+_{\geq i'}\bfP$.
This completes the proof that (ii) implies (iii).

It remains to show (ii).
Before doing so, pick any linear extension of the partial order on $\bbZ^2$ where $(i,j)\geq (i',j')$ if and only if $(i,j)$ is weakly northwest of $(i',j')$.
We will show (ii) by induction on the positions $(i,j)$ satisfying the hypotheses of the lemma using the above order.

Suppose $(i,j)$ is contained in a big chute $C$ during step $i$ of $X^+\bfP$.
Let $(i,j_1)$ denote the NE corner of $C$, so $j\leq j_1< j_0$ and $(i,j_1)$ has a black checker in $\bfP$.
Our goal is to show that $j_1=j$.

If instead $j_1> j$, then by (i) $(i,j_1)$ is contained in a big ladder $L'$ during step $j_1$ of $Y^+\bfP$ (note that no such $L'$ can exist for the base case of the induction since its SE corner would be southeast of $(i,j)$, so in this case $j_1 = j$).
Write $(i_1,j_1)\in P_y$ denote the SW corner of $L'$.
We cannot have $i_1 = i$ since otherwise $(i,j_1)$ would still have a black checker in $Y^+_{\geq j_1}\bfP$ in contradiction to (i), so $i_1 > i$.

By the induction hypothesis, $(i_1,j_1)$ satisfies both (ii) and, hence, (iii).
In particular, $(i+1,j_1)$ has a red checker in $X^+_{\geq i+1}\bfP$.
On the other hand, we know $(i+1,j_1)$ has no checker in $X^+_{\geq i+1}\bfP$, since this is the SE corner of the big chute $C$ during step $i$ of $X^+\bfP$.
This gives us the desired contradiction, so we must have $j_1=j$.
This completes the induction for (ii).
\end{proof}

\end{document}
